\chapter{Lattices and Fixed Points}
\label{ch:ch10-lattices-fixed-points}

Iteration can loop forever when its update rule has no notion of progress.
A dataflow solver that grows and shrinks the same set on alternate visits never settles.
A constant propagator that flips a variable between two values on each pass never reports.
The problem of this chapter is to give iteration a mathematical frame in which termination and safety follow from the shape of the value domain.
We build that frame from ordered sets with well-behaved combinations, and we show that every analysis from the previous chapter is an instance of it.
The frame also predicts exactly which analyses terminate quickly, which need help on loops, and which answers a client may trust.
Section~\ref{sec:ch10-order} introduces the order, Section~\ref{sec:ch10-transfer} adds monotone transfer functions, and Section~\ref{sec:ch10-worked} computes a fixed point by hand on the graph in Figure~\ref{fig:ch10-loop}.

\section{Partial orders and lattices}
\label{sec:ch10-order}

A \term{partial order} is a set equipped with a relation $\sqsubseteq$ that is reflexive, antisymmetric, and transitive.
Reflexive means every element sits below itself, antisymmetric means two elements below each other are equal, and transitive means steps compose.
Program facts fit naturally, but an analysis must specify its order explicitly.
Under a concretization-based precision order, moving up admits more behaviors; under the reversed order, moving down does so.
The extremal solution sought by a solver depends on that choice.
This reversal of intuition causes most beginner errors with lattices, so we keep information order and safety direction explicit throughout.

A \term{lattice} is a partial order in which every pair of elements has a greatest lower bound and a least upper bound.
The greatest lower bound is the \term{meet}, written $\meet$.
The least upper bound is the \term{join}, written $\join$.
Their meanings depend on the order: under subset inclusion they are intersection and union, while reversing the order exchanges them.
A finite lattice also has a \term{top} element above everything and a \term{bottom} element below everything.
Top and bottom name order-theoretic extrema, not universal meanings such as unknown or contradiction.
In the flat meet-based examples below, top is the unvisited identity and bottom represents conflicting evidence.
Table~\ref{tab:ch10-lattices} lists the small lattices used in this book with their orders, meets, and heights.

Consider the constant lattice for one variable, with top above every integer constant and every constant above bottom.
Two equal constants meet to themselves, because both sides agree.
Two different constants meet to bottom, because no single constant describes both paths.
Join behaves dually, sending distinct constants to top; in this reversed precision order it is meet, not join, that conservatively combines execution paths.
The \term{height} of this lattice is three, since the longest chain from top to bottom passes through one constant.
Height matters because it bounds how many times a block fact can change during iteration.
A height-three scalar cell changes strictly at most twice; a product containing many cells can have a much longer chain.
Counting elements in a chain, a subset lattice over $U$ facts has height $U+1$.
Algorithm~\ref{alg:ch10-fixedpoint} exploits exactly this bound.

\begin{irlisting}
; Loop that strains an analysis without a lattice argument
; B1 initializes i, B2 tests, B3 steps, B4 exits
B1: i = 0
B1: jmp B2
B2: br i < 10, B3, B4
B3: i = i + 1
B3: jmp B2
B4: use i
\end{irlisting}

The loop listing shows why finiteness alone does not guarantee speed.
The fact \code{i = 0} holds on entry to the first iteration and fails on every later iteration.
An analysis that tracks exact values would visit $B2$ once per trip around the loop.
An analysis over the constant lattice instead meets $0$ with $1$, then with bottom, and stops after two changes.
The answer \enquote{bottom} at the loop head means that no single constant describes \code{i} there.
That answer is weak but safe, and it arrives in bounded time.
A client that needs the trip count must use a richer domain, such as intervals from the next chapter, and pay the corresponding height cost.
The lattice makes the trade between precision and cost explicit before anyone runs the solver.

\section{Monotone transfer and fixed points}
\label{sec:ch10-transfer}

A \term{transfer function} maps incoming lattice values to outgoing lattice values for one block.
It is \term{monotone} when $x \sqsubseteq y$ implies $f(x) \sqsubseteq f(y)$, so more precise inputs yield more precise outputs.
It is \term{distributive} when it preserves meets exactly, a stronger property that some frameworks need for maximal precision.
Every gen-kill function from the previous chapter is monotone, and most are distributive over their set lattices.
Constant folding transfer is monotone but not distributive, which is why constant propagation can miss facts that a hypothetical path-by-path execution would find.
Monotonicity is the property the solver actually needs, and distributivity is the property a precision argument might additionally want.
We require the first everywhere and invoke the second only when we claim strongest-fixed-point results.

A \term{fixed point} of the whole system is an assignment of one lattice value to each block that satisfies every equation simultaneously.
The Knaster--Tarski theorem gives least and greatest fixed points of a monotone map over a complete lattice.
Finite-height iteration reaches the appropriate extreme; on infinite-height lattices, ordinary countable Kleene iteration needs additional continuity assumptions to reach a fixed point at its limit.
We start every block at bottom for a least-fixed-point computation that grows, or at top for a greatest-fixed-point computation that shrinks.
We then apply transfers and meets until nothing moves.
Each application can only move values monotonically along chains of bounded height, so the process stops.
The limit satisfies the equations by construction, and minimality or maximality follows from starting at the extreme and moving monotonically.
Safonov treats this discipline as part of what makes an optimizing compiler worthy of trust, see~\cite{safonov-2010-trustworthy}.
A \term{widening} operator accelerates convergence on lattices of infinite or embarrassing height, such as intervals over machine integers.
Widening jumps to a coarser value after a few iterations, for example from $[0,1]$ to $[0, +\infty)$, and thereby forces termination.
It sacrifices precision deliberately, and a companion \term{narrowing} pass may recover some of it afterward.
This chapter uses widening only on loop heads, where unbounded growth originates.
An implementation must document where widening fires, because two compilers that widen differently report different facts from the same equations.
The alternative of iterating to the exact limit is not available when the limit needs unboundedly many steps.
Section~\ref{sec:ch10-implementation} returns to this engineering choice with concrete thresholds.

\begin{algorithm}
\caption{Generic lattice fixed-point iteration with a worklist}\label{alg:ch10-fixedpoint}
\begin{pseudocode}
\Require reachable CFG $G$ with predecessor-free entry, finite-height lattice, monotone transfers, $s=\top$, entry boundary
\Ensure assignment $X[B]$ that satisfies all equations and is extremal from start $s$
\State $X[B] \gets s$ for every block $B$
\State $W \gets$ all blocks in $G$
\While{$W \neq \emptyset$}
\State remove some $B$ from $W$
\State $in \gets \meet_{P \in \Pred(B)} \Tf{P}(X[P])$, or boundary where $B$ has no predecessor
\State $new \gets in$
\If{$new \neq X[B]$}
\State $X[B] \gets new$
\State add $\Succ(B)$ to $W$
\EndIf
\EndWhile
\end{pseudocode}
\end{algorithm}

Algorithm~\ref{alg:ch10-fixedpoint} generalizes the bit-vector solver to any finite-height lattice.
It seeds every block with the extremal start value, then repeatedly folds predecessor contributions through transfers and the meet.
A change at one block schedules its successors, exactly as in the gen-kill case.
Termination follows from finiteness of height: each block value moves monotonically and can move at most $H$ times, where $H$ is the lattice height.
Correctness follows from the same equation check as before, now stated over abstract values rather than bit sets.
The algorithm never inspects what the lattice elements mean.
Backward analyses reverse the dependency direction; ascending analyses use the dual order; infinite-height interval analyses additionally need a convergence strategy.
Morgan describes optimizer architectures that exploit exactly this sharing across passes, see~\cite{morgan-1998-optimizing}.

The seeding direction is part of the specification, not an implementation detail.
A least-fixed-point computation starts at bottom and climbs by joins, suitable for may questions such as reachability.
A greatest-fixed-point computation starts at top and descends by meets, suitable for must questions such as availability.
Swapping the seeds computes the wrong extreme or diverges from the intended safety direction.
A useful sanity check runs the solver twice on a tiny graph, once from each extreme, and confirms that the two answers bracket each other in the order.
When they coincide, the equations determine a unique fixed point and the seeding debate dissolves.
When they differ, the client must state which extreme its safety proof assumes.
\section{Worked example with a sign lattice}
\label{sec:ch10-worked}

We now iterate a sign analysis over the loop in Figure~\ref{fig:ch10-loop} by hand.
Each variable takes a value in the flat sign lattice with elements top, negative, zero, positive, and bottom.
The order places top above every sign, every sign above bottom, and no order between distinct signs.
The meet of equal signs is that sign, the meet of distinct signs is bottom, and the meet with top returns the other side.
Block $B1$ sets \code{i} to zero, block $B3$ adds one, block $B2$ tests the bound, and block $B4$ consumes the result.
Table~\ref{tab:ch10-lattices} summarizes this lattice alongside two others, and Table~\ref{tab:ch10-trace} traces the iteration values at the loop head.

\begin{figure}
\begin{tikzpicture}
\node[basic block] (b1) {B1: i = 0};
\node[basic block, below=1.1cm of b1] (b2) {B2: test i less 10};
\node[basic block, below left=1.1cm and 0.7cm of b2] (b3) {B3: i = i + 1};
\node[basic block, below right=1.1cm and 0.7cm of b2] (b4) {B4: use i};
\draw[cfg edge] (b1) -- (b2);
\draw[cfg edge] (b2) -- (b3);
\draw[cfg edge] (b2) -- (b4);
\draw[back edge] (b3) -- (b2);
\end{tikzpicture}
\caption{A counting loop whose head merges the first-entry value with values carried around the back edge, which forces the sign of \code{i} to bottom.}
\label{fig:ch10-loop}
\end{figure}

\begin{table}
\caption{Three small lattices compared by elements, combination operators, and longest chain height.}
\label{tab:ch10-lattices}
\begin{tabular}{lccc}
\toprule
Lattice & Elements & Meet and join & Height \\
\midrule
Bit-vector subsets & all subsets of facts & union and intersection & facts count + 1 \\
Constant flat & top, constants, bottom & agree or bottom & 3 \\
Sign flat & top, neg zero pos, bottom & agree or bottom & 3 \\
\bottomrule
\end{tabular}
\end{table}

\begin{table}
\caption{Sign of \code{i} at the loop head across Kleene iterates, from the empty start to stabilization at bottom.}
\label{tab:ch10-trace}
\begin{tabular}{lcc}
\toprule
Iterate & Incoming values at B2 & Meet result \\
\midrule
First & zero from B1 only & zero \\
Second & zero and positive & bottom \\
Third & zero and bottom & bottom \\
\bottomrule
\end{tabular}
\end{table}

We seed every block with top, the identity for this meet, and treat the boundary as the value flowing from $B1$ on first entry.
The first visit to $B2$ sees only the entry edge carrying zero, so $X[B2]$ becomes zero.
Processing $B3$ applies the add-one transfer to zero and yields positive on the back edge.
The second visit to $B2$ meets zero from $B1$ with positive from $B3$ and drops to bottom.
The third visit meets zero with bottom and stays at bottom, so no successor is enqueued for that edge again.
The solver stops with $X[B2] = \bot$, $X[B3] = \bot$, and $X[B4] = \bot$.
The hand computation uses two strict head changes and one unchanged visit.
Height bounds strict changes, not visits: several predecessors can schedule a block even when their combined contribution leaves its value unchanged.
A bit-vector problem on the same graph would instead stabilize after at most one change per fact per block.
The \enquote{loss} from zero to bottom is the solver discovering that the loop head genuinely sees two different signs.

Starting from bottom instead would immediately lose the first-iteration zero fact in this meet-based solver.
Bottom means conflicting signs in the chosen reversed order, not an unvisited predecessor.
An implementation must record the lattice order and initialization convention rather than assigning a universal meaning to the words top and bottom.
Section~\ref{sec:ch10-transfer} states the rule, and the exercises ask you to exhibit a system where the two extremes differ.

\section{Implementation notes and cost}
\label{sec:ch10-implementation}

A lattice implementation needs only three operations plus iteration: meet, join, and equality, with transfer functions supplied per analysis.
A flat lattice stores one tag plus one payload, so all three operations are constant time.
A bit-vector lattice stores one bit per fact, so operations cost one pass over machine words.
An interval lattice stores two bounds, so meet is max of lows with min of highs and join is the reverse, each constant time.
The generic solver calls these operations without knowing which lattice it serves.
This separation keeps the trusted core small: the worklist logic is reviewed once, while each domain supplies its own small, testable operators.
Appel keeps the solver and the domain in the same explicit correspondence between theory and code, see~\cite{appel-1998-modern}.

Let $N$ be the block count, $E$ the edge count, $H$ the lattice height, and $\Delta$ the maximum in-degree.
Each block value changes strictly at most $H-1$ times, giving $O(N+EH)$ worklist visits with duplicate-free queue membership.
Each visit may scan all predecessors, so the displayed algorithm uses $O((N+EH)(\Delta+1))$ lattice operations and predecessor-transfer evaluations.
Storage is $O(N)$ lattice elements in addition to the CFG; the size of an element and the cost of a transfer must also be counted.
A scalar flat lattice has $H=3$, but a product of $V$ flat cells has height $2V+1$.
A subset lattice over $U$ facts has height $U+1$.
Intervals over mathematical integers have infinite ascending chains; widening supplies convergence rather than a universal fixed visit bound.
The C standard permits implementations to define overflow, trap, and rounding behavior within stated latitude, so interval bounds must model those choices explicitly, see~\cite{wg14-2011-c11}.
The abstract-interpretation framework of Cousot and Cousot relates fixed-point construction to sound approximation; its assumptions matter when moving beyond finite-height domains~\cite{cousot-1977-abstract}.

\section{Exercises}
\label{sec:ch10-exercises}

\begin{enumerate}
\item \label{ex:ch10-height-proof} Prove that each stored block value in Algorithm~\ref{alg:ch10-fixedpoint} changes strictly at most $H-1$ times on a lattice of height $H$.
Explain why this does not bound worklist visits by $H$, and derive a whole-graph visit bound by counting initial insertions and successor notifications.
Identify the two strict changes and the unchanged visit in Table~\ref{tab:ch10-trace}.
\item \label{ex:ch10-two-extremes} Construct a two-block CFG with a monotone but non-distributive transfer where least and greatest fixed points differ.
Give the lattice, the transfer table, both fixed points, and the safety direction each one serves.
Explain which client should use which extreme, using one concrete optimization decision that flips between them.
\item \label{ex:ch10-widen-loop} Implement the sign analysis from Section~\ref{sec:ch10-worked} over Figure~\ref{fig:ch10-loop} with widening disabled.
Extend the loop bound from $10$ to $1000$ and count head visits to stabilization.
Then replace signs with intervals, rerun without widening, and report how the visit count grows with the bound.
Add widening after two changes and report the new count with its precision loss.
\end{enumerate}

\section{Further reading}
\label{sec:ch10-reading}

Appel connects fixed-point iteration to working optimizer code without hiding the termination argument, see~\cite{appel-1998-modern}.
Morgan shows how shared solver infrastructure serves many optimization passes, see~\cite{morgan-1998-optimizing}.
Safonov treats monotone frameworks as part of a compiler correctness case, see~\cite{safonov-2010-trustworthy}.
The C language definition bounds what value domains may assume about arithmetic and conversion, see~\cite{wg14-2011-c11}.
